\chapter{Haar Measure and Weyl Integration}

\begin{center}\itshape
A random rotation should not favour one orientation over another.
A random unitary matrix should keep its law when multiplied by a fixed
unitary matrix. These requirements determine the probability measure.
\end{center}

Choosing each matrix entry independently does not give a unitary matrix.
Choosing angles independently does not generally give a uniform rotation.
We need a rule that respects the group itself. On a compact matrix group,
Haar probability is that rule: multiplication by a fixed element leaves
the distribution unchanged.

The eigenvalues of a matrix chosen this way are not independent points
on the circle. Their joint density contains a factor for every pair,
$|e^{i\theta_j}-e^{i\theta_k}|^2$. It vanishes when the pair meets and
suppresses small separations. From that finite-dimensional density we
will derive the limiting pair law
\[
 1-\left(\frac{\sin\pi u}{\pi u}\right)^2.
\]
At $u=0$ its value is zero, understood by continuity. Away from zero it
need not be small: repulsion does not forbid close pairs.

On 6 April 1972, Montgomery and Dyson met at the Institute for Advanced
Study after Chowla introduced them. Montgomery described a pair-correlation
formula suggested by his work on zeta zeros; Dyson recognised the expression
from random matrix theory. The accounts and date are documented in the
sources below. The connection brought together questions that had been
studied for different reasons. The matrix law will be proved here.
The corresponding full law for zeta zeros remains conjectural; the
restricted theorem under the Riemann hypothesis will be identified
separately from it.

Invariant matrix integration predates Haar's general theorem. Hurwitz
used invariant volumes in 1897, and Weyl used them in representation
theory in 1925. Haar's 1933 theorem, discussed in Volume~8, Chapter~13,
placed invariant measure in a broader setting. For matrix groups we
can construct it directly from a volume element at the identity.
Matrix size $n$ is at least one; a group's manifold dimension may be zero.

\section{Invariant volume forms}

A top-degree form also defines a positive measure, even if the form itself
has a sign. On an oriented $m$-dimensional chart, write
$\omega=a(x)\,\dd x^1\wedge\cdots\wedge\dd x^m$ and use the density
$|a(x)|\,\dd\lambda(x)$. The change-of-variables computation of
Proposition~\ref{prop:intwelldef}, with the absolute determinant, shows
that these measures agree on overlapping charts. A partition of unity
glues them to a Borel measure, written $|\omega|$.
This is the measure of the form; it is not a new pointwise differential form.

We need three consequences of change of variables and Fubini from
Volume~7, Theorems~9.8 and~9.6. If $\omega$ is positive for the chosen
orientation, then $\int f\,\dd|\omega|=\int_Mf\omega$ for continuous
compactly supported $f$. For any diffeomorphism $F:N\to M$,
\[
 \int_N(g\circ F)\,\dd|F^*\omega|=\int_Mg\,\dd|\omega|
 \quad(g\ge0\text{ Borel}).
\]
The same identity holds for integrable signed or complex functions by
splitting them into integrable parts. Finally, on a product manifold,
$|\alpha\wedge\beta|=|\alpha|\otimes|\beta|$ for top forms pulled
back from the two factors: the chart density is $|a(x)b(y)|$.
In dimension zero a chart is a point and a top form is a scalar;
these conventions give counting measure multiplied by its absolute value.

Let $G$ have dimension $m$ and Lie algebra $\mathfrak g$.
Left translation $L_g$ has differential $X\mapsto gX$, so
$T_gG=g\mathfrak g$ by Theorem~\ref{thm:closedsubgroup}.
Write $R_g$ for right translation. A choice of volume at $I$ can now
be carried to every tangent space by $L_g$.

\begin{theorem}[Haar measure of a matrix Lie group]\label{thm:haarform}
Choose a nonzero alternating $m$-form $\omega_I$ on $\mathfrak g$.
The rule
\[
 \omega_g(gX_1,\ldots,gX_m)=\omega_I(X_1,\ldots,X_m)
\]
defines a smooth nowhere-zero form on $G$ satisfying $L_h^*\omega=\omega$.
Its measure $\mu=|\omega|$ is a left Haar measure.
Moreover, for every $h\in G$ and Borel set $E$,
\[
 R_h^*\omega=\det(\mathrm{Ad}(h^{-1}))\,\omega,\qquad
 \mu(Eh)=|\det\mathrm{Ad}(h^{-1})|\,\mu(E).
\]
Here $\mathrm{Ad}(h)X=hXh^{-1}$ acts on $\mathfrak g$.
\end{theorem}
\begin{proof}
Choose a basis $X_1,\ldots,X_m$ of $\mathfrak g$ and extend its dual
coordinate forms to linear functionals $\ell_i$ on $M_n$.
The forms $\theta^i_g(v)=\ell_i(g^{-1}v)$ are smooth on $GL_n$.
Restricted to $G$, they form the dual frame to $gX_1,\ldots,gX_m$.
Thus
$\omega=c\,\theta^1\wedge\cdots\wedge\theta^m$,
where $c=\omega_I(X_1,\ldots,X_m)\ne0$.
This proves smoothness and nonvanishing, and also supplies an orientation
on every component. For $m=0$ the empty wedge is one and $\omega=c$
on the discrete group, so the same construction gives $|c|$ times counting
measure.

The differential of $L_h$ sends $gX_i$ to $hgX_i$.
Evaluating the definition on this basis proves $L_h^*\omega=\omega$.
Change of variables then gives $\mu(hE)=\mu(E)$.
In charts the measure has a continuous strictly positive density.
A compact set is covered by finitely many small chart neighbourhoods
with bounded densities and bounded coordinate volumes, so has finite
measure. Every nonempty open set contains a chart ball of positive
measure. Local Lebesgue regularity and a countable atlas give regularity
of the measure. These are the Haar measure properties.

For right translation, its differential sends $gX_i$ to
$gX_ih=gh(h^{-1}X_ih)$. Therefore
\[
 (R_h^*\omega)_g(gX_1,\ldots,gX_m)
 =\omega_I(\mathrm{Ad}(h^{-1})X_1,\ldots,
             \mathrm{Ad}(h^{-1})X_m).
\]
Alternating top forms multiply by the determinant of the applied linear
map, proving the form identity. Taking its absolute measure and applying
change of variables to $\mathbf1_{Eh}$ gives the stated right-translation
factor. In dimension zero this determinant is the empty determinant one.
\end{proof}

With the convention of Volume~8, the modular function is
$\delta(h)=|\det\mathrm{Ad}(h^{-1})|$, so $\mu(Eh)=\delta(h)\mu(E)$.
For $GL_n(\R)$ it is one: conjugation on $M_n$ has determinant
$\det(h)^{-n}\det(h)^n=1$.
For affine maps $x\mapsto ax+b$ it is $1/a$, as
Exercise~\ref{ex:axbmodular} will compute. Thus left-invariant volume
need not be right-invariant volume.

\begin{corollary}\label{cor:compactunimodular}
On a compact matrix Lie group, $\delta(h)=1$ for every $h$, so its Haar
measure is also right invariant. If the group is compact and connected,
then its left-invariant volume form itself satisfies $R_h^*\omega=\omega$.
\end{corollary}
\begin{proof}
The map $h\mapsto\det\mathrm{Ad}(h^{-1})$ is a continuous homomorphism
into $\R^\times$. Its image is compact. An element of absolute value
different from one would have unbounded positive or negative integer
powers, so every image value is $1$ or $-1$.
Their absolute values are one, proving the measure statement.
For a connected group the image is connected and contains one, so it
consists only of one, proving the stronger form statement.
\end{proof}

On a compact group the constructed measure has finite positive total
mass; divide by it to obtain a probability. The next argument shows why
no other invariant probability can disagree with it.

\begin{theorem}[uniqueness]\label{thm:haarcompact}
Suppose a compact metrizable group $G$ carries a right-invariant Borel
probability $\mu$. Every finite left-invariant Borel measure $\nu$ then
satisfies $\nu=\nu(G)\mu$.
In particular, a compact matrix Lie group has exactly one left-invariant
Borel probability. It is the normalised measure above, is right invariant,
and satisfies $\mu(E^{-1})=\mu(E)$ for every Borel $E$.
\end{theorem}
\begin{proof}
For continuous $f$, apply Fubini to $f(yx)$ with measures $\nu(x)$ and
$\mu(y)$. Left invariance of $\nu$ makes its integral in $x$ equal to
$\int f\,\dd\nu$ for every $y$. Right invariance of $\mu$ makes its
integral in $y$ equal to $\int f\,\dd\mu$ for every $x$.
Consequently
\[
 \int f\,\dd\nu
 =\int_G\int_Gf(yx)\,\dd\nu(x)\,\dd\mu(y)
 =\nu(G)\int f\,\dd\mu.
\]
Finite Borel measures on a compact metric space are determined by these
continuous-function integrals. Explicitly, for a nonempty compact $K$, the continuous
functions $\max(0,1-jd(x,K))$ decrease to $\mathbf1_K$.
Dominated convergence gives equality on compact sets, and Borel regularity
gives equality on all Borel sets, as in Volume~8, Theorem~13.3.

For compact matrix groups, the preceding corollary makes the constructed
probability right invariant, so the uniqueness result applies to it.
The measure $\check\mu(E)=\mu(E^{-1})$ is another left-invariant
probability: $(hE)^{-1}=E^{-1}h^{-1}$ and $\mu$ is right invariant.
Uniqueness gives $\check\mu=\mu$.
\end{proof}

\begin{corollary}\label{cor:pushhaar}
A continuous surjective homomorphism between compact matrix Lie groups
sends Haar probability to Haar probability.
\end{corollary}
\begin{proof}
For $\varphi:G\to H$, define $\nu(E)=\mu(\varphi^{-1}(E))$.
This is a Borel probability. If $h=\varphi(g)$, then
$\varphi^{-1}(hE)=g\varphi^{-1}(E)$, so $\nu(hE)=\nu(E)$.
The uniqueness theorem on $H$ identifies $\nu$ as its Haar probability.
\end{proof}

\begin{example}[$SU(2)$ and $SO(3)$]\label{ex:su2haar}
The first column identifies $SU(2)$ with
$S^3=\{(a,b)\in\C^2:|a|^2+|b|^2=1\}$.
Left multiplication by $g\in SU(2)$ acts on that column by the unitary
map $z\mapsto gz$, hence by an orthogonal map of $\R^4$.
For a standard real Gaussian vector $Z\in\R^4$, polar coordinates show
that $Z/|Z|$ has uniform spherical probability, as in
Exercise~\ref{ex:spherevolume}. The zero vector has probability zero.
The Gaussian law and the spherical probability are invariant under all
orthogonal maps, so this spherical law is left invariant on $SU(2)$.
Uniqueness makes it Haar probability.

Push it through the spin covering of Theorem~\ref{thm:spin}.
A uniform unit quaternion therefore produces a Haar random spatial
rotation. This is the basis of Shoemake's random-orientation construction.
It does not mean choosing the rotation angle uniformly: the angle density
will be computed in Exercise~\ref{ex:su2weyl}.
\end{example}

\section{Sampling a random unitary matrix}

A complex Gaussian vector already has the symmetry we want. A standard
complex Gaussian is $(X+iY)/\sqrt2$ for independent real standard normals,
with density $\pi^{-1}e^{-|z|^2}$ on $\C$.
An independent $n$-tuple has density $\pi^{-n}e^{-|z|^2}$ on $\C^n$.
Unitary maps preserve both the norm and real Lebesgue measure, so they
preserve this law. Orthonormalising independent such columns will turn
that symmetry into a random unitary matrix.

\begin{lemma}\label{lem:qr}
Each invertible complex matrix $Z$ has a unique factorisation $Z=QR$,
where $Q$ is unitary and $R$ is upper triangular with strictly positive
real diagonal. The map $Z\mapsto Q$ is continuous on $GL_n(\C)$.
A matrix with independent standard complex Gaussian entries is invertible
with probability one.
\end{lemma}
\begin{proof}
Apply Gram--Schmidt to the columns. At step $k$, subtract the projections
on the earlier orthonormal columns and divide the remaining vector by its
positive norm. Linear independence makes that norm nonzero.
This gives $Z=QR$ with the required diagonal and gives continuous formulas
for all columns of $Q$ on $GL_n$.
If $QR=Q'R'$, then $Q'^{-1}Q=R'R^{-1}$ is both unitary and upper
triangular with positive diagonal. Its first column has just its first
entry; norm one and orthogonality force the other first-row entries to
vanish. Repeat on the lower block. The matrix is diagonal with unit-modulus
positive entries, hence the identity. This proves uniqueness.

For Gaussian columns, condition on the first $k-1$ columns.
Their complex span is a proper real linear subspace of $\C^n$ when
$k\le n$, and has Lebesgue measure zero. The independent $k$th column
has a density, so it lies in that span with probability zero.
Fubini and the finitely many steps prove almost-sure independence.
\end{proof}

\begin{theorem}[Gaussian sampling of Haar measure]\label{thm:qrhaar}
Let $Z$ have independent standard complex Gaussian entries. Its positive-
diagonal factorisation $Z=QR$ gives a matrix $Q$ with Haar law on $U(n)$.
\end{theorem}
\begin{proof}
For fixed $g\in U(n)$, multiplying every column by $g$ preserves its
Gaussian law and the columns' independence. Thus $gZ$ has the same law
as $Z$. Its positive-diagonal factorisation is $(gQ)R$ by uniqueness,
so $Q(gZ)=gQ(Z)$. The continuous map $Q$ is Borel on the probability-one
invertible set. Its law is a left-invariant Borel probability on $U(n)$,
which uniqueness identifies as Haar probability.
\end{proof}

\begin{remark}
A numerical QR routine may choose other phases for the diagonal of $R$.
Without a specified convention, the uniqueness argument above does not
identify the law of its $Q$. Some conventions introduce bias; they do not
all necessarily do so. For any invertible QR factorisation, set
$\Lambda=\operatorname{diag}(r_{jj}/|r_{jj}|)$.
Then $(Q\Lambda)(\Lambda^*R)$ has positive diagonal, so $Q\Lambda$
is the factor proved Haar above. Mezzadri's account explains this phase
correction. The real-Gaussian version, with positive diagonal, gives Haar
probability on $O(n)$ by the same invariance proof.
\end{remark}

\section{Weyl's integration formula}

For two eigenvalues, the distance factor is
$|e^{i\theta_2}-e^{i\theta_1}|^2=4\sin^2((\theta_2-\theta_1)/2)$.
The square will come from a complex coordinate, or two real coordinates,
in a change of variables.

A class function is unchanged by conjugation:
$f(gUg^{-1})=f(U)$. Trace, determinant and the number of eigenvalues in
an arc are examples. By Theorem~\ref{thm:torus}, a class function is
specified by its values on the diagonal torus $T$, symmetric in those
entries. Put $t_j=e^{i\theta_j}$ and use the Vandermonde product
\[
 V(t)=\prod_{j<k}(t_k-t_j)=\det(t_k^{j-1})_{1\le j,k\le n}
\]
from Volume~1, Proposition~7.4. For $n=1$ the empty product is one.

We can order the principal eigenangles measurably, even when some eigenvalues
coincide. For $0\le a<2\pi$, the set of matrices with at least $j$
eigenvalues in the closed arc $\{e^{it}:0\le t\le a\}$ is closed.
Indeed, along a convergent sequence of such matrices, select those $j$
roots and then the remaining roots. Compactness of the circle gives a
subsequence on which all roots converge. Continuity of characteristic
polynomial coefficients shows that the limiting roots, with their
multiplicities, are the roots of the limiting matrix, and the selected
ones remain in the closed arc. For increasing principal angles
$\theta_1\le\cdots\le\theta_n$, this closed set is exactly
$\{\theta_j\le a\}$. Thus each angle, and every Borel expression in a
finite list of them, is Borel measurable. Randomly permuting the list
will remove this particular ordering when we describe a joint density.

\begin{theorem}[Weyl, 1925]\label{thm:weyl}
For every bounded Borel class function $f:U(n)\to\C$,
\[
 \int_{U(n)}f\,\dd\mu
 =\frac1{n!}\int_{[0,2\pi]^n}
 f(\operatorname{diag}(e^{i\theta_1},\ldots,e^{i\theta_n}))
 \prod_{j<k}|e^{i\theta_k}-e^{i\theta_j}|^2
 \frac{\dd\theta_1\cdots\dd\theta_n}{(2\pi)^n}.
\]
Equivalently, a Haar unitary matrix's eigenvalues, given an independent
uniform random ordering, have density $|V(t)|^2/n!$ relative to uniform
probability on $T$.
\end{theorem}

To remove the arbitrary phases of eigenvectors, record only their ordered
eigenlines. We can do this by a Hermitian matrix whose eigenvalues are
fixed at $1,\ldots,n$. Changing its eigenvectors changes the lines but
leaves those eigenvalues fixed.
Let $\mathrm{Herm}_n$ be the real vector space of Hermitian matrices,
with inner product $\langle A,B\rangle=\operatorname{tr}(AB)$ and an
oriented Euclidean volume form $\mathrm{vol}$.
Conjugation $c_g(H)=gHg^{-1}$ is a linear isometry on this space.
Its determinant is continuously $1$ or $-1$ as $g$ varies over connected
$U(n)$, and equals one at $I$. Thus $c_g^*\mathrm{vol}=\mathrm{vol}$.

\begin{lemma}\label{lem:orbit}
For $D=\operatorname{diag}(1,\ldots,n)$, the orbit
$\mathcal O=\{gDg^{-1}:g\in U(n)\}$ is a compact submanifold of
$\mathrm{Herm}_n$ of dimension $n^2-n$.
Its tangent space at $D$ consists of the zero-diagonal Hermitian matrices.
There is a smooth nowhere-zero $(n^2-n)$-form $\sigma$ on $\mathcal O$
with $c_g^*\sigma=\sigma$ for every $g\in U(n)$.
\end{lemma}
\begin{proof}
Set $p(H)=(\operatorname{tr}H,\ldots,\operatorname{tr}H^n)$.
The first $n$ power sums of eigenvalues determine the characteristic
polynomial by Newton's identities
\[
 ke_k=\sum_{i=1}^k(-1)^{i-1}e_{k-i}p_i,\qquad e_0=1.
\]
These follow by differentiating
$E(x)=\prod_j(1+\lambda_jx)=\sum_ke_kx^k$ and comparing coefficients
in $xE'(x)=E(x)\sum_{i\ge1}(-1)^{i-1}p_ix^i$.
The spectral theorem therefore gives $p^{-1}(p(D))=\mathcal O$.
The orbit is compact as the continuous image of compact $U(n)$.

Differentiation and cyclicity of trace give
$dp_H(K)=(k\operatorname{tr}(H^{k-1}K))_{1\le k\le n}$.
At $H=gDg^{-1}$, its values on $gE_{jj}g^{-1}$ are the vectors
$(kj^{k-1})_k$. The matrix with these columns is a Vandermonde matrix
times an invertible diagonal matrix, so they are independent.
Thus $p(D)$ is a regular value. The regular value theorem gives the
claimed dimension and tangent space $\ker dp_H$.
At $D$, the same Vandermonde invertibility says that the kernel condition
is exactly $K_{jj}=0$ for every $j$.

The normal vectors $N_k(H)=kH^{k-1}$ are the gradients of $p_k$.
They are independent and span the orthogonal complement of the tangent
space. Define
\[
 \sigma_H(K_1,\ldots,K_{n^2-n})
 =\mathrm{vol}(N_1(H),\ldots,N_n(H),K_1,\ldots,K_{n^2-n}).
\]
This is smooth and nonzero on every tangent basis.
Since $N_k(c_gH)=c_gN_k(H)$ and $c_g$ preserves the ambient volume,
it preserves $\sigma$. For $n=1$ the orbit is a point and this construction
is a nonzero zero-form, consistent with the stated dimension.
\end{proof}

The polynomials $P_j(x)=\prod_{k\ne j}(x-k)/(j-k)$ equal one at $j$
and zero at the other fixed eigenvalues. Consequently $P_j(H)$ is the
orthogonal projection onto $H$'s $j$th eigenline.
The smooth map
\[
 \Phi:\mathcal O\times T\longrightarrow U(n),\qquad
 \Phi(H,t)=\sum_jt_jP_j(H)
\]
therefore satisfies $\Phi(gDg^{-1},t)=gtg^{-1}$.
Both domain and target have dimension $n^2$.
Let $\omega$ be a Haar form on $U(n)$, normalised so its absolute measure
is Haar probability, and let
$\dd\theta=\dd\theta_1\wedge\cdots\wedge\dd\theta_n$ on $T$.
Although angles are defined only modulo $2\pi$, their differentials agree
across coordinate changes and are globally defined.

\begin{lemma}[the Jacobian]\label{lem:jacobian}
A nonzero constant $C$ satisfies
$\Phi^*\omega=C|V(t)|^2\,\sigma\wedge\dd\theta$ on $\mathcal O\times T$.
\end{lemma}
\begin{proof}
Write $\Phi^*\omega=J(H,t)\,\sigma\wedge\dd\theta$.
Conjugation on $U(n)$ preserves $\omega$, since it is a left translation
followed by a right translation and the group is compact and connected.
The identity $\Phi(c_gH,t)=c_g\Phi(H,t)$ and invariance of $\sigma$
give $J(c_gH,t)=J(H,t)$. It is enough to compute at $H=D$.

A tangent vector there is $(K,it\Theta)$, with $K$ Hermitian of zero
diagonal and $\Theta$ real diagonal. Put $X_{jj}=0$ and
$X_{jk}=K_{jk}/(k-j)$ for $j\ne k$.
Then $X^*=-X$ and $[X,D]=K$.
Differentiate the curve $(e^{sX}De^{-sX},te^{is\Theta})$ under $\Phi$.
Its image has derivative
\[
 Xt-tX+it\Theta=t(t^{-1}Xt-X+i\Theta).
\]
Left invariance lets us remove the factor $t$ when evaluating $\omega$.
Use the complex coordinates $K_{jk}$, $j<k$, and real diagonal coordinates
$\Theta_{jj}$ on the source tangent space.
On the target algebra, use its upper-diagonal complex entries and diagonal
entries divided by $i$. The derivative fixes each $\Theta_{jj}$ and
multiplies each $K_{jk}$ by
\[
 z_{jk}=\frac{t_k/t_j-1}{k-j}.
\]
Multiplication by a complex number $z$ has real determinant $|z|^2$.
The full determinant is therefore
\[
 \prod_{j<k}\frac{|t_k-t_j|^2}{(k-j)^2}.
\]
The fixed ratios of our coordinate volume forms contribute a nonzero
constant independent of $t$. This proves the formula, including its zeros
when eigenvalues coincide.
\end{proof}

Call a diagonal matrix regular when its entries are distinct, and call a
unitary matrix regular when its eigenvalues are distinct.
Write $T_{\mathrm{reg}}$ and $U(n)_{\mathrm{reg}}$ for these sets,
and use the subscript $\mathrm{sing}$ for their complements.
Distinct labellings of the eigenlines explain the factor $n!$.

\begin{lemma}\label{lem:sheets}
\begin{enumerate}
\item[(a)] The inverse image of $U(n)_{\mathrm{reg}}$ under $\Phi$ is
$\mathcal O\times T_{\mathrm{reg}}$. Each regular unitary matrix has
exactly $n!$ preimages.
\item[(b)] If a local diffeomorphism $F:X\to Y$ has exactly $N$ points
in each fibre, for a finite integer $N$, then every $y\in Y$ has an
open neighbourhood evenly covered by $N$ open sets, each mapped
diffeomorphically onto that neighbourhood.
\end{enumerate}
\end{lemma}
\begin{proof}
(a) If $\Phi(H,t)=gtg^{-1}$, its eigenvalues are precisely the entries
of $t$, proving the inverse-image assertion.
For a regular matrix, assign the labels $1,\ldots,n$ to its distinct
eigenlines in any order. This determines $H$ uniquely, with label $j$
as its eigenvalue on the assigned line, and determines $t_j$ as the
unitary eigenvalue on that line. Each of the $n!$ assignments works and
no others do.
(b) For $N=0$ the source is empty and the assertion is vacuous.
Otherwise write $F^{-1}(y)=\{x_1,\ldots,x_N\}$ and choose disjoint
local inverse neighbourhoods at these points.
Intersect their image neighbourhoods to obtain an open $W$ about $y$.
Every point of $W$ already has one preimage in each branch.
It cannot have an additional preimage outside them, since its fibre has
exactly $N$ points. Restricting those branches over $W$ therefore gives
the whole inverse image and the required disjoint covering.
\end{proof}

\begin{proof}[Proof of Theorem~\ref{thm:weyl}]
\emph{Remove the singular matrices.}
The Jacobian lemma says that $\Phi$ is critical exactly on
$\mathcal O\times T_{\mathrm{sing}}$, and its image there is precisely
$U(n)_{\mathrm{sing}}$.
Sard's theorem in equal dimensions, Theorem~\ref{thm:sard}, applied in
countably many charts, makes this image null in every target chart.
Haar measure has a continuous density there, so its measure is zero.
The set $T_{\mathrm{sing}}$ is a finite union of diagonal subtori
$t_j=t_k$ and has zero torus volume. For $n=1$ both singular sets are empty.

\emph{Count every regular preimage.}
On the regular set the Jacobian is nonzero, so $\Phi$ is a local
diffeomorphism with exactly $n!$ sheets by Lemma~\ref{lem:sheets}.
Its image is $U(n)_{\mathrm{reg}}$, which is therefore open.
For a Borel set $E$ inside one evenly covered neighbourhood, change of
variables on each sheet gives
\[
 \int_{\mathcal O\times T_{\mathrm{reg}}}
       \mathbf1_E\circ\Phi\,\dd|\Phi^*\omega|=n!\mu(E).
\]
A countable base gives a countable cover by such neighbourhoods.
Cut any Borel set into disjoint pieces, each inside one of them, and add.
Then pass from indicators to bounded Borel functions.
The singular sets just removed contribute zero, so
\[
 \int_{\mathcal O\times T}f\circ\Phi\,\dd|\Phi^*\omega|
 =n!\int_{U(n)}f\,\dd\mu.
\]
This measure identity uses the absolute Jacobian, hence $|C|$, irrespective
of the orientation signs of the chosen forms.

\emph{Integrate out the eigenlines and normalise.}
For a class function, $f(\Phi(H,t))=f(t)$.
The Jacobian and product-measure rules give
\[
 n!\int f\,\dd\mu
 =|C|\left(\int_{\mathcal O}\dd|\sigma|\right)
       \int_Tf(t)|V(t)|^2\,\dd\theta.
\]
The coefficient is positive and independent of $f$.
Expand both Vandermonde determinants:
\[
 |V(t)|^2=\sum_{\pi,\rho\in S_n}
 \operatorname{sgn}(\pi)\operatorname{sgn}(\rho)
 \prod_ke^{i(\pi(k)-\rho(k))\theta_k}.
\]
Circle integration kills a term unless $\pi(k)=\rho(k)$ for every $k$.
Exactly $n!$ terms survive, each with integral $(2\pi)^n$.
Since $\mu(U(n))=1$, setting $f=1$ determines the coefficient as
$1/(n!(2\pi)^n)$, proving the formula.
For randomly ordered eigenvalues, average any test function over all
permutations. This symmetric test is a class function; the density is
symmetric too. The formula therefore gives the asserted joint law.
\end{proof}

\begin{think}
Each pair of eigenvalues contributed one complex coordinate to the
Jacobian calculation. Why does that produce a squared distance?
For real symmetric, complex Hermitian and quaternionic Hermitian matrix
ensembles, the corresponding off-diagonal coordinates have one, two and
four real dimensions. Dyson's quoted threefold classification leads to
exponents $1,2,4$ in those ensembles. Haar orthogonal matrices have additional
spectral constraints and are not the real-symmetric ensemble.
\end{think}

\begin{remark}
The general Weyl formula for a compact connected Lie group is a quoted
extension here. If $T$ is a maximal torus, $\mathfrak t$ its algebra and
$W=N(T)/T$ its Weyl group, it reads
\[
 \int_G f\,\dd\mu=\frac1{|W|}\int_T f(t)
 \det(\mathrm{Ad}(t^{-1})-I)_{\mathfrak g/\mathfrak t}\,\dd t
\]
for class functions, with normalised Haar measure on $T$.
The general maximal-torus theory needed for this statement has not been
proved here; see Bump in the sources.
For $U(n)$ we have proved the required torus facts and $|W|=n!$.
The determinant on the quotient is $\prod_{j\ne k}(1-t_k/t_j)=|V(t)|^2$,
matching our Jacobian computation. Weyl used such integration formulas
with character orthogonality to study representations; Volume~10 takes
up representation theory.
\end{remark}

\section{Correlations, and Dyson's formula}

The Vandermonde density is a product of two determinants.
Integrating such a product is easier than expanding its many terms afresh
each time. The following identity performs that expansion once.

\begin{lemma}[Andr\'eief, 1886]\label{lem:andreief}
Let $\nu$ be a finite measure on a measurable space $X$, and let
$\phi_1,\ldots,\phi_n$, $\psi_1,\ldots,\psi_n$ be bounded measurable
complex functions. Then
\[
 \frac1{n!}\int_{X^n}\det(\phi_j(x_k))\det(\psi_j(x_k))
       \,\dd\nu(x_1)\cdots\dd\nu(x_n)
 =\det\left(\int_X\phi_j\psi_k\,\dd\nu\right)_{j,k\le n}.
\]
\end{lemma}
\begin{proof}
Put $M_{ab}=\int\phi_a\psi_b\,\dd\nu$.
Expand both determinants and apply Fubini, which is valid for these
bounded functions and finite measures. The unnormalised integral becomes
\[
 \sum_{\pi,\rho\in S_n}\operatorname{sgn}(\pi)\operatorname{sgn}(\rho)
 \prod_k M_{\pi(k),\rho(k)}.
\]
For fixed $\pi$, relabel $a=\pi(k)$ and put $\tau=\rho\pi^{-1}$.
The product is $\prod_aM_{a,\tau(a)}$, and the sign is
$\operatorname{sgn}\tau$. Summing over $\rho$ therefore gives
$\det M$. There are $n!$ choices of $\pi$, giving the required factor.
\end{proof}

\begin{theorem}[Heine's identity]\label{thm:heine}
For a bounded Borel function $f$ on the unit circle, write
\[
 \widehat f(m)=\int_0^{2\pi}f(e^{i\theta})e^{-im\theta}
                         \frac{\dd\theta}{2\pi}.
\]
If $U$ is Haar on $U(n)$ with eigenvalues $\lambda_1,\ldots,\lambda_n$,
then
\[
 \mathbb E\prod_{j=1}^nf(\lambda_j)
 =\det(\widehat f(j-k))_{0\le j,k<n}.
\]
\end{theorem}
\begin{proof}
The product is a bounded Borel class function, so apply Weyl's formula.
Write $|V|^2$ as the product of $\det(e^{ij\theta_k})$ and
$\det(e^{-ij\theta_k})$, with rows $j=0,\ldots,n-1$.
Multiplying the first determinant's $k$th column by $f(e^{i\theta_k})$
puts the whole integrand in Andr\'eief's form.
Its matrix entries are
$\int f(e^{i\theta})e^{i(j-k)\theta}\,\dd\theta/(2\pi)
=\widehat f(k-j)$. This is the transpose of the stated matrix and has
the same determinant.
\end{proof}

A matrix whose entries depend only on the difference of their indices
is called Toeplitz. Heine's identity turns an eigenvalue question into
one of these determinants. Choosing $f$ to be the indicator of the
complement of an arc, for example, gives the probability that the arc
contains no eigenvalue. Exercise~\ref{ex:gap} makes the calculation.

For joint counts, introduce the kernel
\[
 K_n(\theta,\varphi)=\sum_{a=0}^{n-1}e^{ia(\theta-\varphi)}.
\]
Orthogonality of circle exponentials gives
\[
 \int K_n(\theta,\psi)K_n(\psi,\varphi)
           \frac{\dd\psi}{2\pi}=K_n(\theta,\varphi),\qquad
 K_n(\theta,\theta)=n.
\]
The finite geometric sum also gives
$|K_n(\theta,\varphi)|=
|\sin(n(\theta-\varphi)/2)/\sin((\theta-\varphi)/2)|$.
At coincident angles modulo $2\pi$ this quotient means its continuous
value $n$, not a division by zero.

\begin{lemma}[Gaudin]\label{lem:gaudin}
Let $\nu$ be finite and $K$ bounded and measurable, with
\[
 \int K(x,y)K(y,z)\,\dd\nu(y)=K(x,z),\qquad
 \int K(x,x)\,\dd\nu(x)=c.
\]
For every integer $m\ge1$,
\[
 \int\det(K(x_a,x_b))_{a,b\le m}\,\dd\nu(x_m)
 =(c-m+1)\det(K(x_a,x_b))_{a,b\le m-1}.
\]
For $m=1$ the determinant on the right is the empty determinant one.
\end{lemma}
\begin{proof}
Expand the determinant as a sum over $\pi\in S_m$.
Terms with $\pi(m)=m$ have a factor $K(x_m,x_m)$ integrating to $c$.
They leave exactly $c$ times the $(m-1)$-by-$(m-1)$ determinant.

In every other term, the cycle containing $m$ has a predecessor $j<m$
and a successor $l<m$. The two factors involving $x_m$ are
$K(x_j,x_m)K(x_m,x_l)$, which integrate to $K(x_j,x_l)$.
Removing $m$ from that cycle reverses its permutation sign.
Each permutation on $m-1$ symbols arises exactly $m-1$ times: insert
$m$ after any of its symbols in its cycle.
These terms contribute $-(m-1)$ times the smaller determinant.
Adding the two contributions proves the formula. When $m=1$ only the
first contribution exists, giving $c$ as required.
\end{proof}

\begin{theorem}[correlation functions of random unitary matrices]\label{thm:correlations}
Let the eigenangles of a Haar $U\in U(n)$ be
$\theta_1,\ldots,\theta_n\in[0,2\pi)$.
For $1\le k\le n$ and every bounded Borel function $g$ on $[0,2\pi)^k$,
\[
 \mathbb E\sum_{j_1,\ldots,j_k\ \mathrm{distinct}}
       g(\theta_{j_1},\ldots,\theta_{j_k})
 =\int g(\varphi_1,\ldots,\varphi_k)
       \det(K_n(\varphi_a,\varphi_b))_{a,b\le k}
       \frac{\dd\varphi_1\cdots\dd\varphi_k}{(2\pi)^k}.
\]
The sum is over ordered tuples of distinct indices and is independent
of the ordering chosen for the eigenangles.
\end{theorem}
\begin{proof}
For $A=(e^{ij\theta_l})_{0\le j<n,\ 1\le l\le n}$,
the Gram determinant identity gives
\[
 |V|^2=|\det A|^2=\det(A^*A)
       =\det(K_n(\theta_a,\theta_b))_{a,b\le n}.
\]
The last matrix may be transposed relative to $A^*A$; transposition does
not change its determinant.
Use the randomly ordered joint law from Weyl's theorem.
Its symmetry makes all $n!/(n-k)!$ terms of the sum contribute the same
integral. The expectation is therefore
\[
 \frac1{(n-k)!}\int g(\theta_1,\ldots,\theta_k)
       \det(K_n(\theta_a,\theta_b))_{a,b\le n}
       \frac{\dd\theta_1\cdots\dd\theta_n}{(2\pi)^n}.
\]
The kernel satisfies Gaudin's assumptions with $c=n$, by the orthogonality
identities above. Integrating out $\theta_n,\ldots,\theta_{k+1}$
produces factors $1,2,\ldots,n-k$.
They cancel the factorial, giving the result. For $k=n$ there are no
integrations to perform and the factorial is $0!=1$.
\end{proof}

For one angle the intensity is $n$, relative to normalised circle measure.
Thus an arc of length $\ell$ contains $n\ell/(2\pi)$ eigenvalues on
average. For two angles the intensity is
$n^2-|K_n(\theta,\varphi)|^2$, vanishing at a collision.
Its vanishing makes close pairs less likely, rather than impossible.
To see the local law as $n$ grows, measure separations in units of the
mean spacing $2\pi/n$.

\begin{corollary}[Dyson's pair correlation]\label{cor:dyson}
For real $x$, let $\langle x\rangle\in(-\pi,\pi]$ be its representative
modulo $2\pi$. For every continuous compactly supported function $h$ on
$\R$,
\[
 \lim_{n\to\infty}\frac1n\mathbb E\sum_{j\ne k}
 h\left(\frac n{2\pi}\langle\theta_j-\theta_k\rangle\right)
 =\int_\R h(u)\left[1-\left(\frac{\sin\pi u}{\pi u}\right)^2\right]\dd u.
\]
For $n$ independent uniform circle points, the corresponding limit is
$\int_\R h(u)\,\dd u$.
\end{corollary}
\begin{proof}
For $n\ge2$, apply the two-point correlation formula.
Its integrand depends only on the circular difference
$x=\langle\varphi_1-\varphi_2\rangle$.
For fixed $\varphi_2$, this difference preserves uniform circle measure.
The expectation divided by $n$ becomes
\[
 \frac1n\int_{-\pi}^{\pi}h\left(\frac{nx}{2\pi}\right)
 \left[n^2-\frac{\sin^2(nx/2)}{\sin^2(x/2)}\right]\frac{\dd x}{2\pi}.
\]
Set $x=2\pi u/n$. The result is
\[
 \int_{-n/2}^{n/2}h(u)
 \left[1-\frac{\sin^2\pi u}{n^2\sin^2(\pi u/n)}\right]\dd u.
\]
The bracket lies in $[0,1]$ because $|K_n|\le n$.
At $u=0$ it has continuous value zero.
For fixed $u\ne0$, $n\sin(\pi u/n)\to\pi u$, giving the limiting
bracket; at zero the same limit holds by the convention just stated.
For all sufficiently large $n$ the integration interval contains the
support of $h$. The absolute integrand is bounded by $|h(u)|$,
which is integrable. Dominated convergence applies even when $h$ has
negative or complex values and proves the limit.

For independent points, each ordered pair has uniform circular difference.
The same substitution gives
$(n-1)n^{-1}\int_{-n/2}^{n/2}h(u)\,\dd u$, tending to $\int h$.
No assertion for a two-point correlation at $n=1$ is needed for either limit.
\end{proof}

\begin{figure}[ht]
\centering
\fig{ch08-cue}
\caption{Left: the $24$ eigenvalues of a Haar random matrix of $U(24)$, sampled as in Theorem~\ref{thm:qrhaar}. Middle: $24$ independent uniform points. The eigenvalues are spread more
evenly; the independent points form clumps and leave gaps. Right: the pair correlation $1-(\frac{\sin\pi u}{\pi u})^2$ of Corollary~\ref{cor:dyson}, against the constant $1$ of independent points;
near $u=0$ it behaves like $\frac{\pi^2}3u^2$.}
\label{fig:cue}
\end{figure}

\begin{beauty}
Under the Riemann hypothesis, write the positive-ordinate nontrivial zeros
of $\zeta$ as $1/2+i\gamma_j$, listed with multiplicity.
Rescale large positive ordinates by
$\widetilde\gamma=\gamma\log(\gamma/(2\pi e))/(2\pi)$.
The Riemann--von Mangoldt formula from Volume~6, Chapter~6, gives
$N(T)=\widetilde T+O(\log T)$, so this scale has mean density one
as height increases.
The full pair-correlation conjecture predicts, for $a\le b$,
\[
 \frac1{N(T)}\#\{(j,k):j\ne k,\ 0<\gamma_j,\gamma_k\le T,
       \ \widetilde\gamma_j-\widetilde\gamma_k\in[a,b]\}
 \longrightarrow\int_a^b
 \left[1-\left(\frac{\sin\pi u}{\pi u}\right)^2\right]\dd u.
\]
Distinct indices, not necessarily distinct numerical ordinates, are counted;
this keeps multiplicities in the question.
Montgomery's 1973 theorem, assuming the Riemann hypothesis, establishes a
weighted, restricted test-function version with rapidly decreasing smooth test functions whose Fourier transforms have
compact support in $(-1,1)$. It is not a proof of the full assertion above.
Odlyzko's computations of high zeros tested the prediction numerically;
those comparisons are evidence and illustrations, not proofs.

The Hilbert--P\'olya proposal seeks a spectral construction whose
self-adjointness would imply the Riemann hypothesis.
The required construction is not known. Self-adjointness alone would
not force the sine-kernel spacing law: different self-adjoint operators
can have different spectral statistics. An explanation of that law
would need additional structure. Volume~12 studies the random-matrix
side of this connection further.
\end{beauty}

\begin{remark}
Dyson's circular unitary ensemble is precisely Haar probability on $U(n)$.
The determinant methods connect it to the Gaussian ensembles studied by
Mehta and Gaudin. Their local sine-kernel limits require further analysis,
including orthogonal polynomials, and are not proved by the circle
calculation here. Those limits are a preview of Volume~12.
\end{remark}

\begin{figure}[ht]
\centering
\fig{v9-ch08-jacobian-square}
\caption{For $j=1,k=2,t_1=1,t_2=i$, the Jacobian multiplies the complex coordinate $K_{12}$ by $i-1$. A unit square of area one becomes a diamond of area two, exactly $|t_2-t_1|^2$. Each off-diagonal pair is a real two-dimensional plane, which produces the square in the Vandermonde weight.}
\label{fig:v9-jacobian-square}
\end{figure}
For an angle set $A$, let $N_A$ count the eigenvalues in it.
Write $\nu=\dd\theta/(2\pi)$ for uniform circle probability and
$p=\nu(A)$. The next picture compares the variance of this count for an arc
of that fraction of the circle. Its formulas are derived from the correlation theorem in
Exercise~\ref{ex:v9-arc-variance}.

\begin{figure}[ht]
\centering
\fig{v9-ch08-count-variance}
\caption{For two points and an arc occupying a fraction $p$ of the circle, independent sampling gives count variance $2p(1-p)$. Haar $U(2)$ eigenvalues give $2p-2p^2-2(\sin(\pi p)/\pi)^2$, the lower curve. The calculation in Exercise~\ref{ex:v9-arc-variance} measures more even counts without claiming that every individual configuration is evenly spaced.}
\label{fig:v9-count-variance}
\end{figure}

\section{What we proved, and what comes next}

The group symmetry determined Haar probability, and Gaussian columns gave
a way to sample it. Passing to eigenvalues changed the volume element:
each complex off-diagonal coordinate contributed a squared pair distance.
That Jacobian explains the Vandermonde weight in Weyl's formula.

The weight also explains a distinction between a uniform average and
independence. Selecting one of the eigenvalues uniformly gives a uniform circle point,
but pairs and larger tuples have determinant correlations.
Andr\'eief's identity and Gaudin's integration step let us compute those
correlations exactly. The rescaled pair law follows by an ordinary limit
and dominated convergence, with no numerical assumption.

The next chapter studies $SL_2(\R)$, which is not compact, acting on the
hyperbolic plane. Its modular subgroup is infinite, yet a fundamental
region has finite invariant area. The measure and the geometry will
again have to be understood together.

\section{Exercises}

\begin{exercise}[$\star$]\label{ex:axbmodular}
For $G=\{\left(\begin{smallmatrix}a&b\\0&1\end{smallmatrix}\right):a>0,b\in\R\}$,
(a) find its Lie algebra and $\det\mathrm{Ad}(h^{-1})$;
(b) compute a left-invariant top form and a right-invariant top form.
Compare with Volume~8, Exercise~13.1, using the convention
$\mu(Eh)=\delta(h)\mu(E)$.
\end{exercise}
\begin{exercise}[$\star\star$]\label{ex:firstcolumn}
Choose $U$ with Haar probability law on $U(n)$.
(a) Show its first column has the law $z/|z|$ for a standard complex
Gaussian vector $z$, hence uniform spherical probability.
(b) For $n\ge2$, prove that $|U_{11}|^2$ has density
$(n-1)(1-x)^{n-2}$ on $[0,1]$. Identify the law separately for $n=1$.
Determine the limiting law of $n|U_{11}|^2$ as $n\to\infty$.
\end{exercise}
\begin{exercise}[$\star\star$]\label{ex:su2weyl}
\emph{The angle of a uniform rotation.}
(a) If a Haar matrix of $SU(2)$ has eigenvalues $e^{\pm i\alpha}$,
$\alpha\in[0,\pi]$, prove its angle density is $2\sin^2\alpha/\pi$.
(b) Show the angle $\phi\in[0,\pi]$ of a Haar spatial rotation has
density $(1-\cos\phi)/\pi$ and mean $\pi/2+2/\pi$, about $126.5^\circ$.
\end{exercise}
\begin{exercise}[$\star\star$]\label{ex:tracemoments}
For Haar $U\in U(n)$ and integer $k\ge1$, prove
$\mathbb E\operatorname{tr}U^k=0$ and
$\mathbb E|\operatorname{tr}U^k|^2=\min(k,n)$.
Compare with the second moment of $\sum_je^{ik\theta_j}$ for $n$
independent uniform circle points. Include $n=1$ in your argument.
\end{exercise}
\begin{exercise}[$\star\star$]\label{ex:gap}
For $0<\alpha<\pi$, prove the probability that Haar $U\in U(n)$ has
no eigenvalue in $|\theta|<\alpha$ equals
$\det(c_{j-k})_{0\le j,k<n}$, where
$c_0=1-\alpha/\pi$ and $c_m=-\sin(m\alpha)/(\pi m)$ for $m\ne0$.
Evaluate it for $n=2$ and check its limits for small $\alpha$ and
$\alpha\to\pi$.
\end{exercise}
\begin{exercise}[$\star$]\label{ex:repulsion}
For Haar $U\in U(2)$, find the joint eigenangle density and the probability
that the smaller circular angular separation is less than $\varepsilon$.
State the boundary cases as well as the formula for $0<\varepsilon<\pi$.
Compare its small-$\varepsilon$ behaviour with independent uniform points.
\end{exercise}

\begin{exercise}[$\star\star$]\label{ex:v9-qr-phase-counterexample}
\emph{A QR convention that changes the group.}
For a Gaussian $Z=QR$ with positive diagonal, put
$D=\operatorname{diag}((\det Q)^{-1},1,\ldots,1)$ and
$Q_0=QD$, $R_0=D^{-1}R$.
Show that $Z=Q_0R_0$ is another unitary--triangular factorisation, but
$Q_0$ is not Haar on $U(n)$. Prove that it is Haar on $SU(n)$ instead.
Show that the phase correction in the sampling remark recovers $Q$.
Include the scalar case $n=1$.
\end{exercise}
\begin{exercise}[$\star\star$]\label{ex:v9-determinant-sampling}
\emph{Separating a scalar phase.}
Let $S$ be Haar on $SU(n)$ and let $W$ be an independent uniform unit
complex number. Prove that $WS$ is Haar on $U(n)$.
How many pairs $(w,s)\in U(1)\times SU(n)$ represent a fixed unitary
matrix this way? Deduce that the determinant of a Haar unitary matrix
is uniform on the circle.
If $F$ is bounded Borel and satisfies $F(wU)=F(U)$ for every $|w|=1$,
prove $\mathbb E[F(U)(\det U)^k]=0$ for every nonzero integer $k$.
\end{exercise}
\begin{exercise}[$\star\star$]\label{ex:v9-arc-variance}
\emph{How much can an eigenvalue count fluctuate?}
Use normalised circle measure $\nu=\dd\theta/(2\pi)$.
For a Borel set of angles $A$, let $N_A$ count the eigenvalues in $A$.
Prove
\[
 \operatorname{Var}N_A=n\nu(A)-\int_{A\times A}|K_n(\theta,\varphi)|^2
                                      \,\dd\nu(\theta)\,\dd\nu(\varphi).
\]
Deduce $0\le\operatorname{Var}N_A\le n\nu(A)$.
For disjoint Borel sets $A,B$, prove
$\operatorname{Cov}(N_A,N_B)=-\int_{A\times B}|K_n|^2\,\dd\nu^2\le0$.
For $n=2$ and an arc occupying fraction $p$ of the circle, compute the
variance explicitly and compare it with two independent points.
Treat $n=1$ separately where the two-point theorem is unavailable.
\end{exercise}
\begin{exercise}[$\star\star$]\label{ex:v9-characteristic-polynomials}
\emph{Averages of characteristic polynomials.}
For arbitrary $z,w\in\C$ and Haar $U\in U(n)$, prove
\[
 \mathbb E\det(I-zU)=1,\qquad
 \mathbb E\left[\det(I-zU)\overline{\det(I-wU)}\right]
       =\sum_{j=0}^n(z\overline w)^j.
\]
Let $e_j(\lambda)$ be the elementary symmetric function of its eigenvalues,
with $e_0=1$, so $\det(I-zU)=\sum_{j=0}^n(-z)^je_j(\lambda)$.
Deduce $\mathbb E[e_j(\lambda)\overline{e_k(\lambda)}]=\delta_{jk}$
and $\mathbb E\det(zI-U)=z^n$.
\end{exercise}


\section{Hints}

\begin{hint}{ex:axbmodular}
Compute the entries of $h^{-1}\dd h$ and $\dd h\,h^{-1}$ and wedge them.
\end{hint}
\begin{hint}{ex:firstcolumn}
Positive-diagonal Gram--Schmidt only normalises the first column.
The squared moduli of independent complex Gaussians are independent
mean-one exponentials.
\end{hint}
\begin{hint}{ex:su2weyl}
The entry $a$ in the spherical coordinates $(a,b)\in S^3$ is uniform
in the unit disc. Its real part is $\cos\alpha$.
\end{hint}
\begin{hint}{ex:tracemoments}
For $n\ge2$, expand the squared trace into a diagonal term $n$ and an
ordered-pair term. Fourier orthogonality counts the remaining pairs.
\end{hint}
\begin{hint}{ex:gap}
Use Heine's identity with the indicator of the complementary arc.
\end{hint}
\begin{hint}{ex:repulsion}
The two-angle density relative to uniform product measure is
$1-\cos(\theta_1-\theta_2)$. Integrate over the circular difference.
\end{hint}

\begin{hint}{ex:v9-qr-phase-counterexample}
The determinant of $Q_0$ is one. For $S\in SU(n)$, check
$Q_0(SZ)=S Q_0(Z)$ using uniqueness of positive-diagonal QR.
\end{hint}
\begin{hint}{ex:v9-determinant-sampling}
Use the surjective homomorphism $(w,s)\mapsto ws$ and choose an $n$th
root of a prescribed determinant. Then integrate the independent phase.
\end{hint}
\begin{hint}{ex:v9-arc-variance}
Expand $N_A^2$ into its one-index and two-distinct-index sums.
For the final calculation, $|K_2|^2=2+2\cos(\theta-\varphi)$.
\end{hint}
\begin{hint}{ex:v9-characteristic-polynomials}
Heine's identity gives a triangular Toeplitz matrix for the first mean.
For the mixed mean it gives a tridiagonal one. Its determinants satisfy
a two-step recurrence, starting with the empty determinant one.
\end{hint}


\section{Solutions}

\begin{solution}{ex:axbmodular}
\idea Conjugation shows how a right translation changes the left volume.
\writeup (a) The algebra consists of
$X=\left(\begin{smallmatrix}x&y\\0&0\end{smallmatrix}\right)$.
Direct multiplication gives
\[
 h^{-1}Xh=\begin{pmatrix}x&(xb+y)/a\\0&0\end{pmatrix}.
\]
Thus on coordinates $(x,y)$ the adjoint map has determinant $1/a$.
(b) The entries above the zero bottom row in $h^{-1}\dd h$ are
$\dd a/a,\dd b/a$. Their wedge is the left-invariant form
$\dd a\wedge\dd b/a^2$.
The corresponding entries of $\dd h\,h^{-1}$ are
$\dd a/a,\dd b-(b/a)\dd a$, whose wedge is
$\dd a\wedge\dd b/a$, a right-invariant form.
Both can also be checked by differentiating the left and right group
products. The theorem gives $\mu(Eh)=a^{-1}\mu(E)$,
matching Volume~8's convention and exhibiting a nonunimodular group.
\end{solution}
\begin{solution}{ex:firstcolumn}
\idea Normalising independent exponential squared lengths gives a simple
ratio whose tail can be integrated exactly.
\writeup (a) In the Gaussian QR sampler, the first column is $z/|z|$.
Unitary invariance of the Gaussian and polar coordinates identify its
spherical law.
(b) Write $E_j=|z_j|^2$. Polar coordinates in $\C$ give
$\Pr(E_j>s)=e^{-s}$ for $s\ge0$, and the $E_j$ are independent.
For $n\ge2$, put $S=E_2+\cdots+E_n$. For $0<x<1$,
\[
 \Pr\left(\frac{E_1}{E_1+S}>x\right)
 =\mathbb E e^{-xS/(1-x)}=(1-x)^{n-1},
\]
since $\mathbb Ee^{-sE_j}=1/(1+s)$.
Differentiating the tail gives the stated density, which integrates to
one and has no endpoint atoms.
For $n=1$, $|U_{11}|^2=1$ identically: its law is the point mass at one,
not the density formula with an exponent $-1$.
For each fixed $s\ge0$ and all sufficiently large $n$,
\[
 \Pr(n|U_{11}|^2>s)=(1-s/n)^{n-1}\longrightarrow e^{-s}.
\]
For $s<0$ the variables are nonnegative. These tails prove convergence
in law to a mean-one exponential random variable.
The argument of $U_{11}$ is the uniform argument of $z_1$, independent
of its modulus; normalising the column does not change that argument.
Combining this independent phase with the radius limit also gives
$\sqrt n U_{11}$ converging in law to a standard complex Gaussian.
Indeed, for any bounded continuous test $F$ on $\C$, averaging
$F(\sqrt s e^{i\theta})$ over uniform $\theta$ gives a bounded continuous
function of $s\ge0$. Apply the squared-radius limit to that function.
\end{solution}
\begin{solution}{ex:su2weyl}
\idea The spherical quaternion law gives a planar marginal, then a change
of variable gives the angle.
\writeup (a) The pair $(a,b)$ has the normalised complex-Gaussian law
on $S^3$. The ratio calculation with two coordinates makes $|a|^2$
uniform on $[0,1]$. Its argument is uniform and independent of the two
Gaussian moduli, hence independent of $|a|$.
Thus $a$ is uniform in the unit disc. Its real part has density
$2\sqrt{1-x^2}/\pi$ on $[-1,1]$, obtained by the length of the vertical
disc slice. Since $\operatorname{tr}U=2\operatorname{Re}a=2\cos\alpha$,
substitution $x=\cos\alpha$ gives $2\sin^2\alpha/\pi$.
(b) The spin covering doubles the angle parameter: if its lift has
angle $\alpha\in[0,\pi]$, the spatial angle in $[0,\pi]$ is
$\phi=\min(2\alpha,2\pi-2\alpha)$.
Both branches contribute equally, and the derivative factor is $1/2$
on each, giving density
$2\sin^2(\phi/2)/\pi=(1-\cos\phi)/\pi$.
The covering sends Haar to Haar by Corollary~\ref{cor:pushhaar}.
Finally integration by parts gives
\[
 \mathbb E\phi=\frac1\pi\int_0^\pi\phi(1-\cos\phi)\,\dd\phi
 =\frac\pi2+\frac2\pi\approx2.207\text{ radians}.
\]
This is about $126.5^\circ$. The density is quadratic near zero, so
small-angle rotations have small probability even under uniform group
sampling.
\end{solution}
\begin{solution}{ex:tracemoments}
\idea Fourier orthogonality turns the pair integral into a count of index
differences.
\writeup The one-point intensity gives mean trace zero for $k\ge1$.
If $n=1$, the trace is one uniform phase raised to $k$, so its squared
modulus is one, proving the assertion directly.
For $n\ge2$, expand
$|\operatorname{tr}U^k|^2=n+\sum_{j\ne l}e^{ik(\theta_j-\theta_l)}$
and use the two-point correlation theorem.
The constant $n^2$ part of its density integrates to zero against this
nonzero Fourier mode. Since
$|K_n|^2=\sum_{a,b=0}^{n-1}e^{i(a-b)(\varphi_1-\varphi_2)}$,
the remaining integral counts pairs with $a-b=-k$.
There are $\max(n-k,0)$, and its minus sign gives
$\mathbb E|\operatorname{tr}U^k|^2=n-\max(n-k,0)=\min(k,n)$.
For independent points all cross terms have mean zero, so the second
moment is $n$ for every $k\ge1$.
In particular the first trace has second moment one for Haar unitaries,
rather than $n$. This second-moment cancellation is one precise meaning
of more regular spacing.
\end{solution}
\begin{solution}{ex:gap}
\idea The product of complementary-arc indicators is exactly the no-point
indicator.
\writeup Set $f(e^{i\theta})=\mathbf1_{[\alpha,2\pi-\alpha]}(\theta)$.
The endpoints do not affect the probability, since the one-point law has
no atoms. Heine's identity gives the determinant of its Fourier coefficients.
The constant coefficient is $1-\alpha/\pi$.
For $m\ne0$, subtract the omitted arc from the full circle to get
\[
 \widehat f(m)=-\frac1{2\pi}\int_{-\alpha}^{\alpha}e^{-im\theta}\,\dd\theta
 =-\frac{\sin(m\alpha)}{\pi m}.
\]
For $n=2$, the probability is
$(1-\alpha/\pi)^2-\sin^2\alpha/\pi^2$.
It tends to zero as $\alpha\to\pi$.
For small $\alpha$ it equals
$1-2\alpha/\pi+(\alpha^2-\sin^2\alpha)/\pi^2
=1-2\alpha/\pi+O(\alpha^4)$.
The linear term is the expected number of points in the omitted arc;
the correction for both eigenvalues being there has order $\alpha^4$,
compared with order $\alpha^2$ for two independent points.
\end{solution}
\begin{solution}{ex:repulsion}
\idea Integrating out the common angle leaves one circular difference.
\writeup Weyl's two-angle density, relative to
$\dd\theta_1\dd\theta_2/(2\pi)^2$, is
$|e^{i\theta_1}-e^{i\theta_2}|^2/2=1-\cos(\theta_1-\theta_2)$.
Thus $x=\langle\theta_1-\theta_2\rangle$ has density
$(1-\cos x)/(2\pi)$ on $(-\pi,\pi]$.
For $0<\varepsilon<\pi$, the desired probability is
\[
 \frac1{2\pi}\int_{-\varepsilon}^{\varepsilon}(1-\cos x)\,\dd x
 =\frac{\varepsilon-\sin\varepsilon}{\pi}.
\]
It is zero for $\varepsilon\le0$ and one for $\varepsilon\ge\pi$;
at $\pi$ the possible excluded antipodal difference has probability zero.
As $\varepsilon\downarrow0$, the probability is
$\varepsilon^3/(6\pi)+O(\varepsilon^5)$.
Independent points instead give $\varepsilon/\pi$ in this range.
Their ratio is asymptotic to $6/\varepsilon^2$; at $\varepsilon=0.1$,
the independent probability is about $600$ times as large. This is a
comparison of probabilities, not an absolute prohibition on close pairs.
\end{solution}

\begin{solution}{ex:v9-qr-phase-counterexample}
\idea Changing a column phase can force determinant one while preserving
a perfectly valid QR factorisation.
\writeup The matrix $D$ is unitary, so $Q_0$ is unitary and $R_0$ is
upper triangular, with $Q_0R_0=QR=Z$.
Also $\det Q_0=(\det Q)(\det D)=1$.
Its law cannot be Haar on $U(n)$: left multiplication by the scalar
$e^{i\pi/n}I$ changes this determinant from one to minus one, and thus
changes its law.
For fixed $S\in SU(n)$, positive-diagonal QR gives $Q(SZ)=SQ(Z)$.
The determinants of $SQ$ and $Q$ agree, so their matrices $D$ agree,
and $Q_0(SZ)=S Q_0(Z)$.
Gaussian left invariance makes the law of $Q_0$ left invariant on $SU(n)$.
It is a Borel probability there and hence Haar by uniqueness.

The first diagonal entry of $R_0$ is $(\det Q)r_{11}$;
the other diagonal entries remain positive.
Its phase matrix is therefore $\Lambda=D^{-1}$, and the corrected
factor is $Q_0\Lambda=Q$.
For $n=1$, $Q_0=1$ and $R_0=Z$; a constant unitary factor is a valid
QR choice but not a uniform circle point. The correction restores its
random phase. This example establishes a possible bias, not a claim
that all numerical conventions have this bias.
\end{solution}
\begin{solution}{ex:v9-determinant-sampling}
\idea A scalar phase supplies the determinant that special unitary matrices
lack.
\writeup The map $U(1)\times SU(n)\to U(n)$, $(w,s)\mapsto ws$,
is a continuous homomorphism, because scalar matrices commute with $s$.
For any $U$, choose $w$ with $w^n=\det U$ and put $s=w^{-1}U$.
Then $s\in SU(n)$, proving surjectivity.
Independent Haar probabilities give Haar probability on the product,
by left invariance and uniqueness. Corollary~\ref{cor:pushhaar} makes
the image $WS$ Haar on $U(n)$.

The equation $w^n=\det U$ has exactly $n$ unit-circle roots, each giving
exactly one $s$. Hence every fibre has $n$ points, including one when
$n=1$. The determinant of $WS$ is $W^n$.
The circle power map is a continuous surjective homomorphism, so $W^n$
is uniform by the same pushforward principle.
For a scalar-invariant $F$, $F(WS)=F(S)$.
Independence and $\det S=1$ therefore give
\[
 \mathbb E[F(WS)(\det(WS))^k]
 =\mathbb EF(S)\,\mathbb EW^{nk}=0
\]
for every nonzero integer $k$, by the integral of a nonconstant circle
Fourier mode. Negative integers cause no denominator problem because
$|W|=1$.
\end{solution}
\begin{solution}{ex:v9-arc-variance}
\idea The two-point density subtracts an explicit positive term from the
independent count.
\writeup Write $p=\nu(A)$. The one-point theorem gives
$\mathbb EN_A=np$.
For $n\ge2$, expand the square into diagonal and distinct-index terms:
\[
 \mathbb EN_A^2=np+n^2p^2
       -\int_{A\times A}|K_n(\theta,\varphi)|^2\,\dd\nu^2.
\]
Subtract $(np)^2$ to obtain the formula.
Variance is nonnegative by its definition, and the subtracted integral
is nonnegative, proving the two bounds.
For disjoint $A,B$, there is no same-index contribution to $N_AN_B$.
The two-point theorem gives
$\mathbb EN_AN_B=n^2\nu(A)\nu(B)-\int_{A\times B}|K_n|^2\,\dd\nu^2$.
Subtracting the product of the means proves the covariance formula.
For $n=1$, $N_A$ is Bernoulli with parameter $p$ and $K_1=1$.
The variance formula becomes $p-p^2$.
For disjoint $A,B$, $N_AN_B=0$ and the covariance is
$-\nu(A)\nu(B)$, again the same formula. Thus no use of correlation
order two is made at $n=1$.

For $n=2$, the kernel identity gives
\[
 \int_{A\times A}|K_2|^2\,\dd\nu^2
 =2p^2+2\left|\int_A e^{i\theta}\,\dd\nu(\theta)\right|^2.
\]
For an arc of fraction $p$, the last modulus is $\sin(\pi p)/\pi$,
by direct integration over a centred arc; rotation does not change it.
Hence
$\operatorname{Var}N_A=2p-2p^2-2(\sin(\pi p)/\pi)^2$.
Two independent points give $2p(1-p)$.
The difference is the nonnegative last term, positive for $0<p<1$;
at $p=0,1$ both counts are deterministic and both variances vanish.
\end{solution}
\begin{solution}{ex:v9-characteristic-polynomials}
\idea The Fourier coefficients leave only three diagonals, so a determinant
recurrence replaces the matrix integral.
\writeup For $f(\lambda)=1-z\lambda$, Heine's Toeplitz matrix is
triangular with all diagonal entries one. Its determinant is one, proving
the first mean.
For the mixed mean, use
\[
 f(\lambda)=(1-z\lambda)(1-\overline w\lambda^{-1})
 =1+z\overline w-z\lambda-\overline w\lambda^{-1}.
\]
On the unit circle, $\lambda^{-1}=\overline\lambda$, so its product
at the eigenvalues is exactly the required mixed determinant.
The Toeplitz matrix has diagonal $1+t$, lower adjacent diagonal $-z$
and upper adjacent diagonal $-\overline w$, where $t=z\overline w$.
Let its size-$m$ determinant be $D_m$. Expansion at an end row gives
\[
 D_m=(1+t)D_{m-1}-tD_{m-2},\qquad D_0=1,\quad D_1=1+t.
\]
Induction verifies $D_m=1+t+\cdots+t^m$.
This is a polynomial identity and includes $t=0$ and $t=1$ without
any division by $1-t$. It proves the mixed mean for all $z,w$.

Compare coefficients in the two-variable polynomial identity.
The left coefficient of $z^j\overline w^k$ is
$(-1)^{j+k}\mathbb E[e_j\overline{e_k}]$;
the right coefficient is one when $j=k$ and zero otherwise.
For equal indices the sign is positive, giving exactly $\delta_{jk}$.
The first mean also shows $\mathbb Ee_j=0$ for $j\ge1$.
Expanding $\det(zI-U)=\sum_{j=0}^n(-1)^je_jz^{n-j}$ then leaves
only $z^n$ after expectation, including at $z=0$.
\end{solution}


\section*{Sources and further reading}
\addcontentsline{toc}{section}{Sources and further reading}

The meeting of Montgomery and Dyson: K.~Sabbagh, \emph{Dr.\ Riemann's Zeros} (Atlantic Books, 2002), pp.~134--136, which quotes Montgomery's account; the date, 6 April 1972, is
given by the Institute for Advanced Study for its conference ``50 years of number theory and random matrix theory'' (2022). H.~L.~Montgomery, ``The pair correlation of zeros of the zeta function'',
in \emph{Analytic Number Theory}, Proc.\ Sympos.\ Pure Math.\ 24 (American Mathematical Society, 1973), 181--193. F.~J.~Dyson, ``Statistical theory of the energy levels of complex systems.~I'',
\emph{J.\ Math.\ Phys.} 3 (1962), 140--156, and ``The threefold way'', same volume, 1199--1215; ``Birds and frogs'', \emph{Notices Amer.\ Math.\ Soc.} 56 (2009), 212--223, where Dyson tells the
story from his side. A.~M.~Odlyzko, ``On the distribution of spacings between zeros of the zeta function'', \emph{Math.\ Comp.} 48 (1987), 273--308.

A.~Hurwitz, ``\"Uber die Erzeugung der Invarianten durch Integration'', \emph{Nachr.\ Ges.\ Wiss.\ G\"ottingen, Math.-Phys.\ Kl.} (1897), 71--90; A.~Haar, ``Der Massbegriff in der Theorie der kontinuierlichen
Gruppen'', \emph{Ann.\ of Math.} 34 (1933), 147--169; H.~Weyl, ``Theorie der Darstellung kontinuierlicher halb-einfacher Gruppen durch lineare Transformationen I--III'', \emph{Math.\ Z.} 23 (1925), 271--309, and
24 (1926), 328--376 and 377--395. J.~von Neumann and E.~Wigner, ``\"Uber das Verhalten von Eigenwerten bei adiabatischen Prozessen'', \emph{Phys.\ Z.} 30 (1929), 467--470. C.~Andr\'eief, ``Note sur une relation
entre les int\'egrales d\'efinies des produits des fonctions'', \emph{M\'em.\ Soc.\ Sci.\ Phys.\ Nat.\ Bordeaux} (3) 2 (1886), 1--14; the volume is often dated 1883, and P.~J.~Forrester, ``Meet Andr\'eief, Bordeaux 1886, and Andreev, Kharkov 1882--1883'', \emph{Random Matrices Theory Appl.} 8 (2019), 1930001, sorts out the date and the author, the Kharkov mathematician K.~A.~Andreev. M.~L.~Mehta and M.~Gaudin, ``On the density of eigenvalues of a random matrix'', \emph{Nuclear Phys.} 18 (1960), 420--427.

F.~Mezzadri, ``How to generate random matrices from the classical compact groups'', \emph{Notices Amer.\ Math.\ Soc.} 54 (2007), 592--604; K.~Shoemake, ``Uniform random rotations'', in \emph{Graphics Gems III}
(Academic Press, 1992), 124--132; P.~Diaconis and M.~Shahshahani, ``On the eigenvalues of random matrices'', \emph{J.\ Appl.\ Probab.} 31A (1994), 49--62.

For more: E.~S.~Meckes, \emph{The Random Matrix Theory of the Classical Compact Groups} (Cambridge University Press, 2019), Chapters~1--3, whose proof of Weyl's formula is close to ours;
D.~Bump, \emph{Lie Groups}, 2nd ed. (Springer, 2013), for Weyl's formula on every compact Lie group; M.~L.~Mehta, \emph{Random Matrices}, 3rd ed. (Elsevier, 2004), for Gaudin's lemma; and the survey of
J.~P.~Keating and N.~C.~Snaith, ``Random matrix theory and $\zeta(1/2+it)$'', \emph{Comm.\ Math.\ Phys.} 214 (2000), 57--89.
